\section{Zadanie 15}
Niech $X_n$ będzie zmienną losową równą sumie oczek w $n$ niezależnych rzutach sprawiedliwą kostką. Pokaż, że dla każdego $k \in \mathbb{N}_2$ zachodzi
\[
\lim_{n \to \infty} P(X_n \text{ jest podzielne przez } k) = \frac{1}{k}.
\]

\subsection*{Rozwiązanie}

Oznaczmy $A_n$ - zmienną losową będącą liczbą oczek w $n$-tym rzucie.

Zauważmy, że:
\[
X_{n+1} \equiv X_n + A_{n+1} \pmod{k},
\]
więc $X_n$ tworzy łańcuch Markowa ze stanami będącymi resztami z dzielenia przez $k$.

Policzymy rozkład stacjonarny tego łańcucha.

Prawdopodobieństwo przejścia z dowolnego stanu $i$ do każdego z $i+1, i+2, \dots, i+6 \pmod{k}$ jest równe $\frac{1}{6}$. Macierz przejścia przyjmie więc postać:
\[
P =
\begin{bmatrix}
\frac{1}{6} & \frac{1}{6} & \frac{1}{6} & \frac{1}{6} & \frac{1}{6} & \frac{1}{6} & 0 & \cdots & 0\\[2ex]
0 & \frac{1}{6} & \frac{1}{6} & \frac{1}{6} & \frac{1}{6} & \frac{1}{6} & \frac{1}{6} & 0 & \cdots\\[2ex]
\vdots & & & & & & & & \vdots\\[2ex]
\frac{1}{6} & \frac{1}{6} & \frac{1}{6} & \frac{1}{6} & \frac{1}{6} & 0 & \cdots & 0 & \frac{1}{6}\\[2ex]
\frac{1}{6} & \frac{1}{6} & \frac{1}{6} & \frac{1}{6} & \frac{1}{6} & \frac{1}{6} & 0 & \cdots & 0
\end{bmatrix}.
\]

Równanie na rozkład stacjonarny $\pi$ przyjmuje postać:
\[
\pi = P \pi,
\]
czyli
\[
\begin{bmatrix}
\pi_0\\[1.5ex]
\pi_1\\[1.5ex]
\vdots\\[1.5ex]
\pi_{k-1}
\end{bmatrix}
=
\begin{bmatrix}
\frac{\pi_1+\pi_2+\pi_3+\pi_4+\pi_5+\pi_6}{6} \\[1.5ex]
\frac{\pi_2+\pi_3+\pi_4+\pi_5+\pi_6+\pi_7}{6} \\[1.5ex]
\vdots\\[1.5ex]
\frac{\pi_0+\pi_1+\pi_2+\pi_3+\pi_4+\pi_5}{6}
\end{bmatrix},
\quad
\text{oraz} \quad
\sum_{i=0}^{k-1} \pi_i = 1.
\]

Jasno widać, że rozwiązaniem tego układu jest:
\[
\pi_i = \frac{1}{k}, \quad i = 0,1,\dots,k-1.
\]

Z własności rozkładu stacjonarnego wynika więc:
\[
\lim_{n \to \infty} P(X_n \equiv i \pmod{k}) = \pi_i =  \frac{1}{k}.
\]

W szczególności:
\[
\lim_{n \to \infty} P(X_n \equiv 0 \pmod{k}) = \pi_0 = \frac{1}{k}.
\]